\chapter{Core-excited spectra XAS/RIXS (menu 37)}
\label{ch:menu37}
\menulabel{37}

\section{Purpose}

The ROCIS module of ORCA computes transition-metal L-edge spectra (manual
section 5.7) and, through the same block, RIXS cross sections; the
CAS-CI/RAS-CI protocol of section 3.13.18 covers core-excited XAS/RIXS in
much the same two-step way.  This menu reads what such a run already
contains: the transition table of its best absorption block, the edge
branching ratio of the spin-orbit-split pair, and the RIXS bookkeeping
(refusal mode, evaluated cross sections, or not requested).  Its writing
modes produce the inputs of the two core-excitation routes instead: the
off-resonance XES of the ROCIS family (a ROCIS input whose RIXS request
makes the engine print the emission tables -- manual section 7.31.4), and
the two-step CAS-CI/RAS-CI protocol of the manual's section 3.13.18 (a
valence SA-CASSCF, then the core orbitals rotated into the active space
for a one-iteration CAS-CI solve).

\section{What you need}

An ORCA output of a ROCIS run: \code{\%rocis} with \code{DoGenROCIS true},
\code{ReferenceMult} matching the ROHF solution, an \code{OrbWin} donor
window on the core shell (e.g.\ the transition metal 2p) and, for the
SOC-corrected tables, \code{rel DoSOC true}.  Optionally answer the second
question (the statistical branching ratio) to get the ratio/stat column.

\section{Walkthrough}

The example reads the [FeCl$_4$]$^{2-}$ fixture: thirty roots per spin
block, \code{DecomposeFosc}, \code{DoSOC} -- fourteen absorption blocks in
the output, of which the SOC-corrected electric-dipole table is the primary
one (2235 state pairs, 935 with non-zero oscillator strength).

\lstinputlisting[style=fbkcode, firstline=54, lastline=70,
  title={The menu-37 example (top: the interaction and the primary-block
  header with the first transitions)}]
  {examples/expected/m37_xas.txt}

\lstinputlisting[style=fbkcode, firstline=142, lastline=150,
  title={The menu-37 example (the elision note, the edge branching and the
  statistical comparison)}]
  {examples/expected/m37_xas.txt}

\section{What you get}

The report \code{<output>.xas.fbk.md}:

\begin{itemize}
  \item the excitation-table summary (root count, multiplicities, lowest
    excitation);
  \item the primary block's transition table (non-zero fosc rows, capped at
    80 with an elision note): initial/final root with their multiplicity
    labels and the energy in eV;
  \item the branching split at the largest energy gap: each cluster's
    count, energy range, fosc-weighted centroid and summed oscillator
    strength, the low/high ratio and -- when a statistical value was given
    -- the ratio/stat column;
  \item the RIXS bookkeeping of the run.
\end{itemize}

Beside the report, the menu writes a plot-ready companion,
\code{<output>.xas.fbk.csv}: every non-zero transition of the primary block
-- without the report's 80-row display cap -- as the state roots with their
labels, the energy in eV and fosc
(\code{i\_root,i\_label,j\_root,j\_label,energy\_eV,fosc}), so the stick
spectrum can be plotted or broadened in any tool
(Chapter~\ref{ch:formats}).

Mode 2 writes a ROCIS XES input instead: the cluster description (structure,
charge, multiplicity, the optional ROHF high-spin preparation, the
\code{XASelems} index, \code{NRoots} and the six-element \code{OrbWin} of
the two donor spaces plus the acceptor, the SOC-corrected-channel and
elastic-line switches, the method keywords) goes to \file{<stem>.xes.inp},
and the checklist to \file{<stem>.xes.inp.fbk.md}: the run's emission
tables carry the ground-state rows (\code{<root>-<mult>A -> 0-<mult>A}) at
the core-emission energies, and
\code{orca\_mapspc <out> XES -x0<lo> -x1<hi> -w<fwhm> -eV -n<npoints>}
renders them to \file{<out>.XES.stk} and \file{<out>.XES.dat} (the
\code{XESSOC} mode for runs with the SOC-corrected channel on).  The
measured recipe is the \file{fecl4\_xes} fixture pair
(\file{fixtures/rocis/}): the generated input runs to normal termination in
about a minute and renders 56 sticks.

Mode 3 writes the two-step CAS-CI core-excited XAS inputs (the manual's
protocol, section 3.13.18): \file{<stem>.casci\_xas.step1.inp} (the valence
SA-CASSCF that writes the \file{.gbw}) and
\file{<stem>.casci\_xas.step2.inp} (MOREAD, the \code{\%scf rotate} of the
core orbitals into the active window, \code{FrozenCore FC\_NONE} and one
CAS-CI iteration over the core-saturated space), with the checklist in
\file{<stem>.casci\_xas.fbk.md}.  The window selection is positional
(measured): the active window is the \code{norb} consecutive orbitals
starting at $N_{\text{inactive}} = (N_e - \mathrm{nel})/2$ -- 42 for the
\file{[FeCl4]2-} probe, 87 for the manual's own Fe(acac)$_3$ example -- and
each core orbital moves to the leading slots of that window through the
five-field \code{\{i, j, 90, 0, 0\}} rotation.  The step-2 run prints the
L-edge transitions (the probe's first at 719.36~eV), and
\code{orca\_mapspc <out> SOCABS ...} renders the SOC-corrected table (785
peaks measured; \code{ABS} for the plain one).

\section{Conventions, cross-checks and boundaries (measured)}

\begin{readbox}
\begin{itemize}
  \item \textbf{The SOC-corrected tables are population weighted}: their
    fosc column header is \code{fosc(D2) (*population)}; the reader keeps
    the printed value and this menu repeats the semantic.  The SOC tables
    list the full state-pair matrix (thousands of rows, mostly zero); the
    report lists the non-zero ones and states both counts.
  \item \textbf{The branching ratio is a data fact}: the two clusters are
    found by the largest energy gap (the natural default for a
    spin-orbit-split edge), and the menu does not name the edges -- the
    assignment depends on the donor/acceptor windows.  The statistical
    references (2:1 for a 2p edge, 3:2 for a 3d edge, from the
    $(2j+1)$ degeneracies) are compared only when the caller passes one;
    deviations indicate electrostatic and spin-orbit effects
    (Thole \& van der Laan 1988; the menu cites the source rather than
    reimplementing its atomic rules).
  \item \textbf{The RIXS window semantics are measured now}: the RIXS
    variant needs the 6-element \code{OrbWin} -- the first four elements
    are the two donor ranges, the last two the acceptor range (manual
    section 7.31.4) -- while a 4-element window gives the engine's
    zero-states refusal.  Measured on the probe: the spin-orbit-split
    core donors \code{6,6,7,8} with a wide acceptor open the RIXS
    channels; putting a valence donor range first gives the zero-states
    refusal, and wide valence donors make the SOC machinery impractical
    (the XES bullets below).  On success the report prints the transition
    counts and the \code{orca\_mapspc ... RIXS} recipe of manual section
    5.7.4.2.
  \item \textbf{The XES route's measured boundaries} (2026-10-02): the
    off-resonance XES is automatic in a RIXS-requested run, and the plain
    (non-SOC) RIXS channel is its carrier -- a root count too small for
    the channel is skipped with the engine's own warning and no XES table
    prints (measured: \code{NRoots} 10 skipped, 30 covered, on the
    \file{fecl4\_xes} probe); the SOC-corrected channel
    (\code{DoRIXSSOC}) additionally stores per-pair transition densities
    (measured past 36~GB transiently at \code{NRoots} 30 on the probe),
    so the writer requests it only on the caller's ask; the six-element
    \code{OrbWin} indices are the target system's own orbital numbers.
  \item \textbf{The KHD XES variant is a documented termination}: with
    \code{DoKHDXESSOC} and its initial/final state lists the engine
    computes the emission intensities and then aborts in its own printout
    (a null stream inside \code{PrintNXESKHDSpectrum}, signal 11 --
    measured for every probed input shape, state list and print level);
    the writer refuses it and names the working routes instead.
  \item \textbf{The two-step CAS-CI protocol's measured boundaries}
    (2026-10-02): the window mechanism is positional
    ($(N_e-\mathrm{nel})/2$; measured on the probe, consistent with the
    manual's own example); the core-orbital indices are the target
    system's own (read from the step-1 output's orbital table); the
    \code{XAS}/\code{XASSOC} mapspc modes do not read these output tables
    (measured refusals) -- \code{SOCABS} and \code{ABS} do; and the
    protocol's first core excitation sits 0.5~eV from the ROCIS result of
    the same system (719.36 vs 718.87~eV measured) -- two independent
    routes to one edge.  The manual's \code{\%rel picturechange /
    FiniteNuc} block made no difference to the probe's CAS-CI L-edge
    energies (measured, kept out of the generated input).
  \item \textbf{Boundaries}: ROCIS's own accuracy statement applies
    (section 5.7.3: several approximations; qualitatively correct
    results).  This menu reads ROCIS outputs, writes the ROCIS XES input
    and writes the two-step CAS-CI XAS inputs; the RAS-CI XES variant
    (section 3.13.19: the same two-step shape with the \code{XESSOC} rel
    keywords and the \code{XASMOs} index) remains a registered
    candidate.
\end{itemize}
\end{readbox}
